% GENERATED FILE. Edit canonical lessons, then run npm run generate:books.
% canonical-source-hash: fnv1a64:5032c7d88f7dcf25
% canonical-lessons: FR-C368-a-condition-que, FR-C368-pourvu-que, FR-C368-a-moins-que, FR-C368-au-cas-ou, FR-C368-sinon

\chapter{À une condition}
\label{ch:fr-a-une-condition}
\hlchaptermodality{368}

\begin{chapteropening}
\textbf{By the end of this chapter:} \emph{I can set a condition: provided that, as long as, unless, in case, or else — with the right mood after each.}

Complete the last lesson of chapter 368: I can set a condition: provided that, as long as, unless, in case, or else — with the right mood after each.
\end{chapteropening}

\section[à condition que]{à condition que — provided that, on condition that}

\label{lesson:FR-C368-a-condition-que}



\begin{quote}
Before the new one: say the French for plausible, then the French for the supposition.
\end{quote}

\subsection*{You'll want to know: à condition que}
\textbf{à condition que} — \textquotedblleft{}provided that, on condition that\textquotedblright{}.

\begin{quote}
Tu peux sortir, à condition que tu sois à la maison avant minuit.
\end{quote}

\begin{quote}
\emph{You can go out, on condition that you are home before midnight.}
\end{quote}

Takes the subjunctive: \emph{que tu sois}. With the same subject, \emph{à condition de} + infinitive.

\begin{cousinweb}
\emph{condition} comes through Old French from Latin \emph{condiciō} \textquotedblleft{}an agreement, terms, a condition\textquotedblright{} (EtymDB), from \emph{condīcere} \textquotedblleft{}to talk over, to agree on\textquotedblright{}, built on \emph{dīcere} \textquotedblleft{}to say\textquotedblright{} (Lewis \& Short).
\end{cousinweb}

\begin{grammarlens}[title={What you've built}]
One more word for telling, explaining and describing at B1.
\end{grammarlens}

\subsection*{Your turn}
\begin{itemize}
\raggedright
  \item \emph{Say it:} \emph{à condition que}
  \item \emph{Say it:} \emph{à condition que}, once more
  \item \emph{Say it:} a sentence of your own with \emph{à condition que}
  \item \emph{Recall:} say the French for plausible, then the French for the supposition, then say \emph{à condition que} again
\end{itemize}

\begin{tcolorbox}[breakable,colback=teal!4,colframe=teal!35!black,title={Before you move on}]
What does à condition que mean? (\textquotedblleft{}Provided that, on condition that\textquotedblright{}.) Say it once more.
\end{tcolorbox}
\section[pourvu que]{pourvu que — as long as, provided that; let us hope that}

\label{lesson:FR-C368-pourvu-que}



\begin{quote}
Before the new one: say the French for the supposition, then the French for provided that.
\end{quote}

\subsection*{You'll want to know: pourvu que}
\textbf{pourvu que} — \textquotedblleft{}as long as, provided that; let us hope that\textquotedblright{}.

\begin{quote}
Nous irons à la plage, pourvu qu'il fasse beau.
\end{quote}

\begin{quote}
\emph{We will go to the beach, as long as the weather is nice.}
\end{quote}

Takes the subjunctive: \emph{qu'il fasse}. At the start of a sentence it is a wish: \emph{Pourvu qu'il vienne !} — let us hope he comes.

\begin{cousinweb}
\emph{pourvu} is the past participle of \emph{pourvoir} \textquotedblleft{}to provide\textquotedblright{}, which comes from Latin \emph{prōvidēre} \textquotedblleft{}to foresee, to provide for\textquotedblright{} (Lewis \& Short), built on \emph{vidēre} \textquotedblleft{}to see\textquotedblright{}.
\end{cousinweb}

\begin{grammarlens}[title={What you've built}]
One more word for telling, explaining and describing at B1.
\end{grammarlens}

\subsection*{Your turn}
\begin{itemize}
\raggedright
  \item \emph{Say it:} \emph{pourvu que}
  \item \emph{Say it:} \emph{pourvu que}, once more
  \item \emph{Say it:} a sentence of your own with \emph{pourvu que}
  \item \emph{Recall:} say the French for the supposition, then the French for provided that, then say \emph{pourvu que} again
\end{itemize}

\begin{tcolorbox}[breakable,colback=teal!4,colframe=teal!35!black,title={Before you move on}]
What does pourvu que mean? (\textquotedblleft{}As long as, provided that; let us hope that\textquotedblright{}.) Say it once more.
\end{tcolorbox}
\section[à moins que]{à moins que — unless}

\label{lesson:FR-C368-a-moins-que}



\begin{quote}
Before the new one: say the French for provided that, then the French for as long as.
\end{quote}

\subsection*{You'll want to know: à moins que}
\textbf{à moins que} — \textquotedblleft{}unless\textquotedblright{}.

\begin{quote}
Je viendrai à pied, à moins qu'il ne pleuve.
\end{quote}

\begin{quote}
\emph{I will come on foot, unless it rains.}
\end{quote}

Takes the subjunctive, often with a \emph{ne} that is not a negative: \emph{à moins qu'il ne pleuve}. With the same subject, \emph{à moins de} + infinitive.

\begin{cousinweb}
\emph{moins} is inherited through Old French \emph{meins} from Latin \emph{minus} \textquotedblleft{}less\textquotedblright{} (EtymDB; Lewis \& Short): at less than this.
\end{cousinweb}

\begin{grammarlens}[title={What you've built}]
One more word for telling, explaining and describing at B1.
\end{grammarlens}

\subsection*{Your turn}
\begin{itemize}
\raggedright
  \item \emph{Say it:} \emph{à moins que}
  \item \emph{Say it:} \emph{à moins que}, once more
  \item \emph{Say it:} a sentence of your own with \emph{à moins que}
  \item \emph{Recall:} say the French for provided that, then the French for as long as, then say \emph{à moins que} again
\end{itemize}

\begin{tcolorbox}[breakable,colback=teal!4,colframe=teal!35!black,title={Before you move on}]
What does à moins que mean? (\textquotedblleft{}Unless\textquotedblright{}.) Say it once more.
\end{tcolorbox}
\section[au cas où]{au cas où — in case}

\label{lesson:FR-C368-au-cas-ou}



\begin{quote}
Before the new one: say the French for as long as, then the French for unless.
\end{quote}

\subsection*{You'll want to know: au cas où}
\textbf{au cas où} — \textquotedblleft{}in case\textquotedblright{}.

\begin{quote}
Je prends un parapluie, au cas où il pleuvrait.
\end{quote}

\begin{quote}
\emph{I am taking an umbrella in case it rains.}
\end{quote}

Takes the conditional, not the subjunctive: \emph{au cas où il pleuvrait}.

\begin{cousinweb}
\emph{cas} comes through Old French \emph{cas} from Latin \emph{cāsus} \textquotedblleft{}a falling, a chance event\textquotedblright{} (EtymDB), from \emph{cadere} \textquotedblleft{}to fall\textquotedblright{} (Lewis \& Short): what may befall.
\end{cousinweb}

\begin{grammarlens}[title={What you've built}]
One more word for telling, explaining and describing at B1.
\end{grammarlens}

\subsection*{Your turn}
\begin{itemize}
\raggedright
  \item \emph{Say it:} \emph{au cas où}
  \item \emph{Say it:} \emph{au cas où}, once more
  \item \emph{Say it:} a sentence of your own with \emph{au cas où}
  \item \emph{Recall:} say the French for as long as, then the French for unless, then say \emph{au cas où} again
\end{itemize}

\begin{tcolorbox}[breakable,colback=teal!4,colframe=teal!35!black,title={Before you move on}]
What does au cas où mean? (\textquotedblleft{}In case\textquotedblright{}.) Say it once more.
\end{tcolorbox}
\section[sinon]{sinon — otherwise, or else}

\label{lesson:FR-C368-sinon}



\begin{quote}
Before the new one: say the French for unless, then the French for in case.
\end{quote}

\subsection*{You'll want to know: sinon}
\textbf{sinon} — \textquotedblleft{}otherwise, or else\textquotedblright{}.

\begin{quote}
Dépêche-toi, sinon nous allons manquer le train.
\end{quote}

\begin{quote}
\emph{Hurry up, or else we will miss the train.}
\end{quote}

It can also mean if not, apart from: \emph{personne, sinon lui}.

\begin{cousinweb}
\emph{sinon} is built from \emph{si} \textquotedblleft{}if\textquotedblright{} and \emph{non} \textquotedblleft{}not\textquotedblright{}, which are inherited from Latin \emph{sī} \textquotedblleft{}if\textquotedblright{} and \emph{nōn} \textquotedblleft{}not\textquotedblright{} (Lewis \& Short): if not.
\end{cousinweb}

\begin{grammarlens}[title={What you've built}]
One more word for telling, explaining and describing at B1.
\end{grammarlens}

\subsection*{Your turn}
\begin{itemize}
\raggedright
  \item \emph{Say it:} \emph{sinon}
  \item \emph{Say it:} \emph{sinon}, once more
  \item \emph{Say it:} a sentence of your own with \emph{sinon}
  \item \emph{Recall:} say the French for unless, then the French for in case, then say \emph{sinon} again
\end{itemize}

\begin{tcolorbox}[breakable,colback=teal!4,colframe=teal!35!black,title={Before you move on}]
What does sinon mean? (\textquotedblleft{}Otherwise, or else\textquotedblright{}.) Say it once more.
\end{tcolorbox}
